\uuid{Akqv}
\titre{Résolution d'équation différentielle du premier ordre}
\auteur{bodin ; croizat ; sacre}
\organisation{exo7}
\licence{CC BY-NC-SA 4.0}
\datecreate{2022-03-08}
\niveau{L1}
\module{Analyse}
\chapitre{Equation différentielle}
\sousChapitre{Résolution d'équation différentielle du premier ordre}
\difficulte{3}
\theme{}
\exo7qcmid{hilisit.equations-differentielles.054}

\contenu{
\begin{question}
\qclassification{225.01}
Soit l'équation différentielle $y' = y + 3x-2$.
Quelles sont les affirmations vraies ?
\begin{answers}
  \good{Une solution particulière est $y(x) = -3x-1$.}
  \bad{Une solution particulière est $y(x) = 3x-2$.}
  \good{La solution vérifiant $y(0)=1$ est $y(x) = 2e^x-3x-1$.}
  \bad{La solution vérifiant $y(0)=1$ est $y(x) = 3e^x +3x-2$.}
\end{answers}
\begin{explanations}
L'équation homogène est $y' = y$, dont les solutions sont les $y_h(x) = Ce^{x}$.
Une solution particulière de l'équation $y' = y + 3x-2$ est $y_p(x) = -3x-1$.
Les solutions générales sont alors $y(x) = y_h(x) + y_p(x)$.
La solutions vérifiant $y(0)=1$ est $y_0(x) = 2e^{x} -3x-1$.
\end{explanations}
\end{question}
}
